\section{Formal Model and the Hybrid Strategy}\label{sec:model}

This section defines the objects shared by the simulator, the oracle and the application, the outcome of an
operation, the metrics and the hybrid strategy. Names are those of the code.

\subsection{Replicas, operations and causality}

\begin{definition}[System]\label{def:system}
A system is a finite set of replicas $R = \{s\} \cup C$, where $s$ is the server and $C = \{c_1, \dots, c_n\}$ are
the clients. Clients exchange messages only with $s$ (star topology). The replicated data is a set of
\emph{subscription records}, each identified by a UUID generated by the replica that creates it. A record has the
fields of Table~\ref{tab:classes}; each field belongs to one \emph{field class}.
\end{definition}

\begin{table}[ht]
\caption{Field classes of a subscription record and the mechanism used for each class by the hybrid strategy S4.}
\label{tab:classes}
\footnotesize
\begin{tabular}{@{}l>{\raggedright\arraybackslash}p{0.34\textwidth}>{\raggedright\arraybackslash}p{0.36\textwidth}@{}}
\toprule
Class & Fields & Mechanism in S4 \\
\midrule
SIMPLE & name, currency, billingCycle, category, isActive, notes & LWW register ordered by HLC \\
CRITICAL & amountCents, nextPaymentDate & multi-value register, conflict surfaced \\
SET & tags & observed-remove set, add wins \\
APPEND & payments & grow-only set keyed by payment id \\
COUNTER & sharedWith & PN-counter \\
RECORD & existence (create, delete) & causal tombstone, edit wins \\
\bottomrule
\end{tabular}
\end{table}

\begin{definition}[Operation]\label{def:op}
Every user intent is one operation $o$ of a kind in $\{\mathit{create}, \mathit{setField}, \mathit{addTag},
\mathit{removeTag}, \mathit{addPayment}, \mathit{increment}, \mathit{delete}\}$ on one record. The operation has an
origin replica $\mathit{origin}(o)$, a sequence number $\mathit{seq}(o)$ among the operations of its origin, a true
time $\tau(o)$ (the simulator's virtual time, independent of the possibly skewed clock of the device) and a version
vector $\mathit{seen}(o) : R \to \mathbb{N}$ counting, per origin, the operations that its replica had already
integrated when $o$ was applied.
\end{definition}

\begin{definition}[Happens-before]\label{def:hb}
For operations $a, b$, $a \to b$ ($a$ happened before $b$) iff $\mathit{seq}(a) \le
\mathit{seen}(b)(\mathit{origin}(a))$, that is, iff $a$ is in the causal history of $b$. Operations $a \neq b$ are
\emph{concurrent}, $a \parallel b$, iff neither $a \to b$ nor $b \to a$.
\end{definition}

Causal histories are prefix-closed per origin, so $\mathit{seen}$ represents them exactly.

\begin{definition}[Strategy]\label{def:strategy}
A strategy is a tuple
\[ S = (\Sigma, \mathit{init}, \mathit{apply}, \mathit{msg}, \mathit{recv}, \mathit{read}, \mathit{conflicts}) \]
where $\Sigma$ is the set of replica states, $\mathit{apply} : \Sigma \times O \to \Sigma$
integrates a local operation, $\mathit{msg} : \Sigma \times R \to (\Sigma \times M) \cup \{\bot\}$ produces the next
message for a peer ($\bot$ when there is nothing to send), $\mathit{recv} : \Sigma \times R \times M \to \Sigma$
integrates a message, $\mathit{read}(\sigma)$ is the visible set of records and $\mathit{conflicts}(\sigma)$ the set
of unresolved conflicts, each with at least two candidate values. A \emph{sync session} between $a$ and $b$
alternates $\mathit{msg}$ and $\mathit{recv}$ in both directions until both sides return $\bot$.
\end{definition}

\begin{definition}[Convergence]\label{def:convergence}
A run \emph{converges} if, after every client has completed two sessions with the server with no new operations,
all replicas have equal $\mathit{read}$ and equal $\mathit{conflicts}$. A strategy is \emph{quiescent} if, after a
session, $\mathit{msg}$ returns $\bot$ on both sides until the next $\mathit{apply}$, and \emph{idempotent} if
receiving the same message twice has the effect of receiving it once.
\end{definition}

All strategies pass a shared property-based conformance suite for these properties and for persistence.

\subsection{Outcome of an operation and metrics}

The oracle judges the final, quiescent state of the server against the log of all operations.

\begin{definition}[Outcome]\label{def:outcome}
Each operation $a$ receives exactly one outcome, checked in this order:
\begin{itemize}
\item \emph{ambiguous}: $a$ is a $\mathit{delete}$ concurrent with an edit of the same record, or an
  $\mathit{addTag}$ concurrent with a $\mathit{removeTag}$ of the same tag (and the other operation of the pair);
  both results are legitimate policy choices, so these operations are reported separately and never counted as lost;
\item \emph{preserved}: the effect of $a$ is visible: the field has the value written by $a$, the tag is present
  (absent after a removal), the payment is present, the record is present (absent after a delete);
\item \emph{superseded}: the effect is not visible, but an operation $b$ with $a \to b$ on the same target
  legitimately replaced it (a delete supersedes every causally earlier operation on its record);
\item \emph{surfaced}: the value of $a$ is not visible but is a candidate of an entry of $\mathit{conflicts}$;
\item \emph{lost}: none of the above.
\end{itemize}
Counter operations are judged together: with unit increments, the number of lost increments of a record is
$|\,\text{expected sum} - \text{actual value}\,|$.
\end{definition}

\begin{definition}[Metrics]\label{def:metrics}
The metrics of a run with operation set $O$ are those of Table~\ref{tab:metrics}.
\end{definition}

\begin{table}[ht]
\caption{Metrics of a run with operation set $O$ (names as in the code).}
\label{tab:metrics}
\footnotesize
\begin{tabular}{@{}l>{\raggedright\arraybackslash}p{0.62\textwidth}@{}}
\toprule
Metric & Definition \\
\midrule
$\mathit{lossRate}$ & lost operations divided by $|O|$ \\
$\mathit{concurrentLossRate}$ & lost operations divided by the operations with at least one concurrent operation on
  the same record \\
$\mathit{surfacedRate}$ & surfaced operations divided by $|O|$ \\
$\mathit{causalViolations}$ & scalar fields whose final value comes only from operations $a$ such that some $b$
  with $a \to b$ wrote the same field (a stale write won) \\
$\mathit{converged}$ & Definition~\ref{def:convergence} \\
$\mathit{payloadBytes}$ & sum of the lengths of all messages, also per session and per operation \\
$\mathit{syncRounds}$ & messages per session \\
$\mathit{syncCpuMs}$ & time spent in $\mathit{msg}$ and $\mathit{recv}$, per session \\
$\mathit{stateBytes}$ & length of the persisted state, per client and for the server \\
$\mathit{metadataRatio}$ & $\mathit{stateBytes}$ divided by the length of the JSON encoding of $\mathit{read}$ \\
$\mathit{loc}$ & non-blank, non-comment lines of code of a strategy \\
\bottomrule
\end{tabular}
\end{table}

\subsection{The strategies}
S1 (\emph{lww}) keeps whole records and orders the last write to a record by the pair (device time, replica id); the
later record replaces the other one entirely, and a delete is a timestamped tombstone. S2 (\emph{vv}) keeps a version
vector per record and a dot per field; concurrent writes of the same field become a conflict with a provisional value,
and a delete concurrent with an edit becomes a conflict on the record's existence. S3 delegates everything to a CRDT
library: one Automerge document, with Yjs as a control implementation, holding a map from record id to record. S4,
described next, is our proposal.

\subsection{The type-aware hybrid strategy (S4)}
The idea of S4 is that each field class gets the cheapest replicated type that keeps its intended semantics
(Table~\ref{tab:classes}), instead of one mechanism for the whole record.
SIMPLE fields are LWW registers stamped with a hybrid logical clock~\cite{kulkarni2014hlc}, a triple (wall time,
counter, replica id) ordered lexicographically. A causally later write always has a larger stamp, even on a device
whose clock runs late, so skew cannot make a stale value win; a concurrent write with the smaller stamp is still lost,
by design, since these fields rarely justify interrupting the user. CRITICAL fields (price, next payment date) are
multi-value registers: concurrent writes are kept as siblings, the one with the largest HLC is shown provisionally and
the others are surfaced as a conflict, resolved by a new write that dominates all siblings. The tag set is an
add-wins observed-remove set without tombstones for removed elements~\cite{bieniusa2012orset,almeida2018delta},
payments form a grow-only set keyed by payment id, and the shared-with counter is a PN-counter with one positive and
one negative entry per replica.

Each record has one causal context (a version vector plus a set of isolated dots) shared by its entries; every local
operation adds one dot, and a delete stores the context it has seen. Under the default policy \emph{editWins} the
record stays visible when its context contains a dot that no delete has seen, that is, a concurrent edit; the
alternative \emph{deleteWins} hides it on any delete. Synchronisation is delta-based: for each peer the replica
remembers which entries changed since the peer's last acknowledgement and sends only those, retransmitting until they
are acknowledged.

\begin{proposition}\label{prop:sec}
S4 is strongly eventually consistent: in every execution in which each delta is eventually delivered, any two
replicas that have received the same set of deltas have equal $\mathit{read}$ and $\mathit{conflicts}$.
\end{proposition}

\begin{proof}
Each building block (HLC-ordered LWW register, multi-value register, add-wins set, grow-only set, PN-counter, causal
context) is a state-based CRDT: its states form a join-semilattice and its merge is the join, hence commutative,
associative and idempotent~\cite{shapiro2011comprehensive,almeida2018delta}. The state of a record is the product of
the states of its entries, and the state of a replica is a map from record ids to record states merged pointwise;
products and pointwise maps of join-semilattices are join-semilattices with the component-wise
join~\cite{shapiro2011crdt}. Every delta $d$ produced by an update satisfies $\sigma \sqcup d = \sigma'$, where
$\sigma'$ is the state after the update~\cite{almeida2018delta}, so a replica's state is the join of its initial
state and the deltas it has received. Since the join is associative, commutative and idempotent, this value does not
depend on the order or multiplicity of delivery, and two replicas with the same set of deltas have the same state.
Finally, $\mathit{read}$ and $\mathit{conflicts}$ are functions of the state alone. The server forwards every delta
it integrates to the other clients and retransmits until acknowledged, which gives eventual delivery in the star
topology whenever every client reconnects. The laws used here are also checked by property-based tests.
\end{proof}
